\documentclass[12pt,a4paper]{article}

\usepackage[a4paper,margin=1in]{geometry}
\usepackage[T1]{fontenc}
\usepackage[utf8]{inputenc}
\usepackage{lmodern}
\usepackage{hyperref}
\usepackage{graphicx}
\usepackage{float}
\usepackage{booktabs}
\usepackage{longtable}
\usepackage{xcolor}
\usepackage{listings}
\usepackage{array}


\hypersetup{
    colorlinks=true,
    linkcolor=blue,
    urlcolor=blue,
    citecolor=blue
}

\lstdefinestyle{cppstyle}{
    language=C++,
    basicstyle=\ttfamily\small,
    keywordstyle=\color{blue},
    commentstyle=\color{teal},
    stringstyle=\color{red!70!black},
    numbers=left,
    numberstyle=\tiny,
    stepnumber=1,
    numbersep=8pt,
    showstringspaces=false,
    breaklines=true,
    frame=single
}

\newcommand{\coursename}{BLG 478E Network Security}
\newcommand{\reporttitle}{Secure Communication Between Alice and Bob Using a KDC}
\newcommand{\studentname}{İhsan Gök}
\newcommand{\studentid}{150210093}
\newcommand{\submissiondate}{19/04/2026}

\newcommand{\ksa}{\texttt{ALICE\_MASTER\_KEY\_32\_BYTES\_LEN\_!!}}
\newcommand{\ksb}{\texttt{BOB\_\_MASTER\_KEY\_\_32\_BYTES\_LEN\_!!}}

\newcommand{\sessionkeyvalue}{\texttt{<paste session key Ks from terminal>}}
\newcommand{\ticketvalue}{\texttt{<paste Bob ticket ciphertext from terminal>}}
\newcommand{\ciphertextvalue}{\texttt{<paste encrypted payload sent to Bob>}}
\newcommand{\plaintextvalue}{\texttt{\studentid\ Network Sec.}}
\newcommand{\hashvalue}{\texttt{a1e625a897566e5514a76d74660a9358b05a8e7ed967fd43aa4bc933973b7c17}}
\newcommand{\tamperedhashvalue}{\texttt{<paste modified hash value used in integrity test>}}

\newcommand{\imageplaceholder}[2]{
    \IfFileExists{#1}{
        \includegraphics[width=\linewidth]{#1}
    }{
        \fbox{
            \parbox[c][0.28\textheight][c]{0.92\linewidth}{
                \centering
                Screenshot placeholder\\[4pt]
                \texttt{#1}\\[8pt]
                #2
            }
        }
    }
}

\newcommand{\reportfigure}[4]{
    \begin{figure}[H]
        \centering
        \imageplaceholder{#1}{#4}
        \caption{#2}
        \label{#3}
    \end{figure}
}

\begin{document}

\begin{titlepage}
    \centering
    {\Large \coursename \par}
    \vspace{1.5cm}
    {\huge \bfseries \reporttitle \par}
    \vspace{1.5cm}
    {\Large Analysis Report \par}
    \vspace{2cm}

    \begin{tabular}{rl}
        Student Name: & \studentname \\
        Student ID: & \studentid \\
        Date: & \submissiondate \\
    \end{tabular}

    \vfill

    This report documents the implementation and network analysis of a secure
    communication system between Alice and Bob using a Key Distribution Center (KDC).
\end{titlepage}

\tableofcontents
\newpage

\section{Introduction}
This project implements a simple secure communication protocol between Alice and Bob
with the assistance of a Key Distribution Center (KDC). The KDC creates a fresh
session key \(K_s\) for each communication request. Alice receives \(K_s\) encrypted
with her long-term master key, while Bob receives the same session key inside a
ticket encrypted with Bob's long-term master key. Alice then uses the session key
to encrypt the message and its SHA-256 hash before sending them to Bob.

The source code is organized into separate programs for the three entities required
by the assignment:

\begin{itemize}
    \item \texttt{kdc.cpp}: KDC server
    \item \texttt{alice.cpp}: Alice client
    \item \texttt{bob.cpp}: Bob server
    \item \texttt{common.cpp} and \texttt{common.h}: shared crypto, encoding, and socket helpers
\end{itemize}

\section{Implementation Overview}
\subsection{Cryptographic Design}
The implementation uses the following primitives:

\begin{itemize}
    \item AES-256-CBC for confidentiality
    \item SHA-256 for message integrity verification
    \item A randomly generated 32-byte session key \(K_s\)
    \item A random 16-byte IV for each AES encryption
\end{itemize}

The IV is prefixed to each ciphertext so the receiver can decrypt the message using
the same IV. The ciphertexts and binary values are encoded as hexadecimal strings
before they are transmitted over the socket connection.

\subsection{Protocol Flow}
The message sequence implemented in the project is:

\begin{enumerate}
    \item Alice sends \texttt{REQ||Alice||Bob} to the KDC.
    \item The KDC generates a fresh session key \(K_s\).
    \item The KDC sends back:
    \[
    \texttt{KDC\_RESP || } E_{K_{A,KDC}}(K_s) \texttt{ || } E_{K_{B,KDC}}(K_s || Alice)
    \]
    \item Alice decrypts \(E_{K_{A,KDC}}(K_s)\) to recover \(K_s\).
    \item Alice computes \(H(M)\) using SHA-256.
    \item Alice sends Bob:
    \[
    \texttt{TICKET || } E_{K_{B,KDC}}(K_s || Alice)
    \]
    followed by
    \[
    \texttt{DATA || } E_{K_s}(M || H(M))
    \]
    \item Bob decrypts the ticket, obtains \(K_s\), decrypts the message, computes
    its own hash, and compares both hashes.
\end{enumerate}

\reportfigure
    {figures/kdc_req.png}
    {Wireshark capture of Alice's request to the KDC on port 5000}
    {fig:kdc-req}
    {Captured packet containing \texttt{REQ||Alice||Bob}.}

\reportfigure
    {figures/kd_resp.png}
    {Wireshark capture of the KDC response carrying encrypted key material}
    {fig:kdc-resp}
    {Captured packet containing \texttt{KDC\_RESP} with the encrypted session key and Bob's ticket.}

\subsection{Source Files}
\begin{longtable}{>{\raggedright\arraybackslash}p{0.24\linewidth} >{\raggedright\arraybackslash}p{0.68\linewidth}}
\toprule
\textbf{File} & \textbf{Responsibility} \\
\midrule
\endhead
\texttt{kdc.cpp} & Accepts Alice's request, generates a fresh session key, encrypts the key for Alice, prepares Bob's ticket, and returns both values. \\
\texttt{alice.cpp} & Connects to the KDC, obtains the session key, hashes the plaintext message, encrypts the payload, and sends the ticket and encrypted payload to Bob. \\
\texttt{bob.cpp} & Accepts Alice's ticket and encrypted data, decrypts them, recomputes the message hash, and prints either \texttt{Verification Successful} or \texttt{Verification Failed}. \\
\texttt{common.cpp/h} & Provides AES, SHA-256, random number generation, hex encoding, and framed socket transmission helpers. \\
\bottomrule
\end{longtable}

\section{Key and Value Definitions}
Table~\ref{tab:keys} lists the long-term keys and important values used by the implementation.

\begin{table}[H]
    \centering
    \begin{tabular}{p{0.28\linewidth} p{0.62\linewidth}}
        \toprule
        \textbf{Item} & \textbf{Value / Description} \\
        \midrule
        Alice identity & \texttt{Alice} \\
        Bob identity & \texttt{Bob} \\
        Alice-KDC master key \(K_{A,KDC}\) & \ksa \\
        Bob-KDC master key \(K_{B,KDC}\) & \ksb \\
        Plaintext message \(M\) & \plaintextvalue \\
        SHA-256 hash \(H(M)\) & \hashvalue \\
        Session key \(K_s\) & \sessionkeyvalue \\
        Bob's encrypted ticket & \ticketvalue \\
        Alice's encrypted payload to Bob & \ciphertextvalue \\
        \bottomrule
    \end{tabular}
    \caption{Main keys and data values used in the implementation}
    \label{tab:keys}
\end{table}

The message text is currently formed in \texttt{alice.cpp} as the student number
concatenated with the string \texttt{" Network Sec."}. Because the message is fixed
in the current code, the SHA-256 hash shown above is deterministic and can be
reported directly.

\section{Baseline Transmission}
To provide a baseline for comparison, the same application-level message should be
observed in Wireshark without applying the AES encryption step. In the plaintext
baseline experiment, the message content should be directly visible in the captured
packet payload, and the string \texttt{150210093 Network Sec.} should appear in
readable form.

This baseline demonstrates why confidentiality is necessary: an observer on the
network can recover both the message and its integrity-related fields if they are
sent without encryption.

\reportfigure
    {figures/baselines.png}
    {Baseline Wireshark capture showing plaintext application data}
    {fig:baseline}
    {Insert a Wireshark screenshot where the plaintext message is visibly highlighted.}

\reportfigure
    {figures/plainmode.png}
    {Alice terminal output during the plaintext baseline run}
    {fig:alice-plaintext}
    {Alice running with \texttt{PLAINTEXT\_MODE=1} for the baseline comparison.}

\noindent\textbf{Observation.}
In the baseline capture, the packet payload is human-readable. This means an
attacker sniffing the traffic can directly see the message and any additional fields
sent alongside it. Such a transmission does not provide confidentiality.

\section{Secured Transmission}
In the secured implementation, the payload sent from Alice to Bob is encrypted with
AES-256-CBC using the fresh session key \(K_s\). Wireshark still captures the
traffic, but the application data is no longer human-readable because it appears as
hex-encoded ciphertext rather than plaintext.

\reportfigure
    {figures/secured-wireshark.png}
    {Wireshark capture of Alice-to-Bob secured transmission}
    {fig:secured}
    {Insert a Wireshark screenshot where the ciphertext is highlighted instead of readable plaintext.}

\noindent\textbf{Observation.}
In the secured capture, the payload does not reveal the original message content.
Although the packet is visible at the network layer, the plaintext message is hidden
because only encrypted data and ticket material are transmitted.

\section{Proof of Implementation}
The terminal outputs of KDC, Alice, and Bob provide direct evidence that the protocol
was implemented and executed correctly. These outputs should clearly show the
inputs and outputs of the encryption and decryption steps.

\subsection{KDC Terminal Output}
The KDC should display the request source and destination identities and the newly
generated session key \(K_s\).

\reportfigure
    {figures/kdc_baselines.png}
    {KDC terminal output during the protocol run}
    {fig:kdc-terminal}
    {Capture the KDC terminal showing the request and generated session key.}

\subsection{Alice Terminal Output}
Alice should display the recovered session key, the ticket received for Bob, the
plaintext message \(M\), and the hash value \(H(M)\).

\reportfigure
    {figures/secured_transmission.png}
    {Alice terminal output during the protocol run}
    {fig:alice-terminal}
    {Capture Alice's terminal showing the session key, ticket, plaintext, and hash.}

\subsection{Bob Terminal Output}
Bob should display the decrypted ticket, recovered session key, sender identity,
received message, received hash, locally calculated hash, and the final verification result.

\reportfigure
    {figures/bob_baselines.png}
    {Bob terminal output during a successful verification}
    {fig:bob-terminal}
    {Capture Bob's terminal showing ticket decryption and successful verification.}

\subsection{Step-by-Step Value Summary}
Table~\ref{tab:proof} is intended to summarize the critical values observed in the
terminal screenshots. Replace the placeholder fields with the actual values from one
complete execution trace.

\begin{table}[H]
    \centering
    \begin{tabular}{p{0.30\linewidth} p{0.60\linewidth}}
        \toprule
        \textbf{Field} & \textbf{Observed Value} \\
        \midrule
        Alice \(\rightarrow\) KDC request & \texttt{REQ||Alice||Bob} \\
        Generated session key \(K_s\) & \sessionkeyvalue \\
        Alice's decrypted \(K_s\) & \sessionkeyvalue \\
        Bob ticket ciphertext & \ticketvalue \\
        Plaintext \(M\) & \plaintextvalue \\
        Computed hash \(H(M)\) & \hashvalue \\
        Encrypted payload sent to Bob & \ciphertextvalue \\
        Bob's calculated hash & \hashvalue \\
        Final result & \texttt{Verification Successful} \\
        \bottomrule
    \end{tabular}
    \caption{Step-by-step protocol values observed during execution}
    \label{tab:proof}
\end{table}

\section{Integrity Test}
To test the integrity mechanism, one bit of the transmitted hash \(H(M)\) should be
modified before Bob receives the encrypted payload. After decryption, Bob will
recompute the hash of the plaintext message and compare it with the tampered value.
Because the two values are different, Bob must print \texttt{Verification Failed}.

This experiment proves that the receiver can detect an unauthorized modification
even when the original message \(M\) itself is left unchanged.

\reportfigure
    {figures/integrity_test.png}
    {Bob terminal output after manually flipping one bit in the transmitted hash}
    {fig:integrity}
    {Insert a screenshot where Bob prints \texttt{Verification Failed}.}

\begin{table}[H]
    \centering
    \begin{tabular}{p{0.30\linewidth} p{0.60\linewidth}}
        \toprule
        \textbf{Field} & \textbf{Value} \\
        \midrule
        Original hash \(H(M)\) & \hashvalue \\
        Tampered hash value & \tamperedhashvalue \\
        Bob's recomputed hash & \hashvalue \\
        Bob's final decision & \texttt{Verification Failed} \\
        \bottomrule
    \end{tabular}
    \caption{Integrity test results after one-bit modification}
\end{table}

\noindent\textbf{Interpretation.}
The integrity check succeeds because Bob computes a fresh SHA-256 hash from the
decrypted plaintext and compares it against the received hash. Any bit flip in the
transmitted hash changes the hexadecimal value, so the comparison fails and Bob
detects the manipulation.

\section{Conclusion}
The implementation satisfies the core assignment requirements by using a KDC to
distribute a fresh session key, AES-256-CBC to provide confidentiality, and SHA-256
to provide integrity verification. The Wireshark comparison shows the difference
between plaintext and encrypted transmission, while the terminal traces show the
full end-to-end protocol execution. Finally, the integrity test confirms that Bob can
detect tampering when the transmitted hash is modified.

\end{document}
